\documentclass{article}
\usepackage{graphicx}
\usepackage{hyperref}
\usepackage{booktabs}
\usepackage{longtable}
\usepackage[utf8]{inputenc}
\usepackage{xeCJK}
\setCJKmainfont{SimSun}
\usepackage{geometry}
\geometry{a4paper, margin=1in}
\title{Modeling Interpretation Brief}
\author{Model Interpreter Agent}
\date{\today}
\begin{document}
\maketitle
\section*{Q4 建模内参 (Modeling Brief)}
\par
\subsection*{1. 变量与参数字典}
\par
\begin{center}
\begin{longtable}{lllll}
\toprule
符号 & 代码变量名 & 物理含义 & 典型值/单位 & 备注 \\
\midrule
$S_t$ & `S\_bc` & 平流层黑碳(BC)累积库存 & ton & 核心状态变量，随发射累积 \\
$\tau$ & `tau\_bc\_years` & 黑碳自然沉降半衰期 & 4.0 Years & 决定了"漏桶"的泄漏速度 \\
$E_{CO2}$ & `E\_CO2\_total\_ton` & 全生命周期碳排放当量 & ton & 包含火箭燃料燃烧与电梯电力消耗 \\
$\chi$ & `chi\_green` & 电力去碳化比例 & 0.0 - 1.0 & 0=煤电, 1=100\%绿电 \\
$N_{safe}$ & `N\_safe\_year` & 臭氧层恢复能力的临界发射频次 & 1000次/年 & 软约束阈值 \\
$EDI$ & `EDI\_proxy` & 环境破坏综合指数 & 无量纲 & 归一化后的多目标加权值 \\
\bottomrule
\end{longtable}
\end{center}
\par
\subsection*{2. 核心建模逻辑}
\par
\subsubsection*{2.1 平流层"漏桶"效应 (The Stratospheric Leaky Bucket)}
不同于对流层污染物能被降水快速清除（几天内），平流层缺乏垂直对流，污染物停留时间极长。我们采用一阶动力学方程描述黑碳的累积与衰减：
\par
\[  S(t+1) = (1 - \delta) \cdot S(t) + u(t)  \]
\par
其中：
*   $S(t)$ 为第 $t$ 天的平流层黑碳存量。
*   $u(t)$ 为当天的黑碳注入量（取决于火箭发射次数）。
*   $\delta = \frac{1}{\tau \times 365}$ 为日衰减率。
\par
**物理意义**：即使单日排放量不大，只要持续注入，且 $\tau$ (4年) 较长，库存量 $S\_t$ 就会持续累积，形成"长尾效应"。这解释了为何即使停止发射，环境影响也会持续数十年。
\par
\subsubsection*{2.2 帕累托权衡 (Pareto Trade-off)}
模型揭示了**建设速度**与**环境破坏**之间的根本矛盾。
*   **快速建设** $\rightarrow$ 需要高频火箭发射 $\rightarrow$ $S\_t$ 峰值突破安全阈值 $\rightarrow$ EDI 指数飙升。
*   **慢速建设** $\rightarrow$ 依赖太空电梯逐步运输 $\rightarrow$ 碳排放受 $\chi$ 控制 $\rightarrow$ 环境风险可控。
\par
\subsection*{3. 可视化图表集总与解读}
\par
\subsubsection*{3.1 黑碳库存的动态演化}
\begin{figure}[h!]
\centering
\includegraphics[width=0.8\textwidth]{outputs/q4/figs/fig_q4_soot_timeseries.png}
\caption{黑碳时间序列}
\end{figure}
> **图 4.1：平流层黑碳库存 ($S\_t$) 的时变轨迹**
> *   **对应需求**：此图即为日级动态轨迹 (`black\_carbon\_trace.png`)。
> *   **X轴**：项目天数；**Y轴**：平流层黑碳存量 (吨)。
> *   **解读**：在项目初期（火箭运输主导阶段），黑碳库存呈快速上升趋势。随着太空电梯（绿色区间）运力占比提升，火箭发射频率降低，$S\_t$ 达到峰值后开始缓慢衰减。
> *   **结论**：必须严格控制初期的火箭发射速率，防止峰值过早突破不可逆阈值。
\par
\subsubsection*{3.2 速度与环境的帕累托权衡}
\begin{figure}[h!]
\centering
\includegraphics[width=0.8\textwidth]{outputs/q4/figs/fig_q4_pareto_edi_cost.png}
\caption{EDI与工期权衡}
\end{figure}
> **图 4.2：完工时间 vs 环境破坏指数 (EDI) 帕累托前沿**
> *   **对应需求**：此图即为时间-环境权衡 (`pareto\_time\_edi.png`)。
> *   **X轴**：项目总工期 (年)；**Y轴**：环境破坏指数 (EDI)。
> *   **趋势**：曲线呈明显的 L 型。试图过度压缩工期（向左移动）会导致 EDI 指数呈指数级恶化。
> *   **拐点**：在约 20-25 年的工期处存在最佳平衡点（Knee Point），此时既能保证效率，又能将环境代价控制在低位。
\par
\subsubsection*{3.3 能源结构的敏感性}
\begin{figure}[h!]
\centering
\includegraphics[width=0.8\textwidth]{outputs/q4/figs/fig_q4_chi_sensitivity.png}
\caption{去碳化敏感性}
\end{figure}
> **图 4.3：电力去碳化比例 ($\chi$) 对总碳排放的影响**
> *   **X轴**：去碳化比例 $\chi$ (0.0 - 1.0)；**Y轴**：总碳排放量。
> *   **解读**：当太空电梯承担主要运力时，其电力来源成为最大碳源。随着 $\chi$ 趋近于 1.0，总排放显著下降。
> *   **关键发现**：若不配套绿电 ($\chi < 0.5$)，太空电梯的总碳排放甚至可能超过高效火箭方案。
\par
\subsubsection*{3.4 环境影响成分拆解}
\begin{figure}[h!]
\centering
\includegraphics[width=0.8\textwidth]{outputs/q4/figs/fig_q4_env_components.png}
\caption{环境成分拆解}
\end{figure}
> **图 4.4：不同方案下的环境破坏成分构成**
> *   **图例**：颜色区分了 CO2、黑碳 (Soot)、臭氧 (O3) 和 建设LCA 的贡献。
> *   **解读**：纯火箭方案的主要环境代价是 **辐射强迫(Soot)** 和 **臭氧损耗**；而太空电梯方案的主要代价是 **电力碳排放**。这验证了"污染转移"假说，并强调了混合运输策略的必要性。
\par
\subsection*{4. 政策建议 (Policy Implications)}
\par
基于上述建模分析，我们提出以下核心建议：
\par
1.  **实施"环境熔断"机制 (Environmental Circuit Breaker)**：
政策不应仅设定静态的"每年发射上限"，而应基于实时监测的 $S\_t$ 库存。当 $S\_t$ 接近临界值（如 $S\_{critical}$）时，必须强制暂停非紧急火箭发射，利用 $\tau$ 周期进行自然清洗。
\par
2.  **绿电强制配套标准**：
仿真结果表明，太空电梯的环保属性高度依赖于能源结构。建议立法规定：**太空电梯项目必须配套建设专属的可再生能源电厂**，确保 $\chi \ge 0.9$。在 $\chi < 0.5$ 的现有电网条件下，盲目建设太空电梯并不能带来显著的气候效益。
\par
3.  **最优工期规划**：
避免盲目追求"最快完工"。图 4.2 显示，将工期从 15 年放宽至 25 年，可使环境破坏指数降低 60\% 以上。应将"环境容忍度"纳入工期规划的约束条件。
\par
---
*Generated by Claude Code - Modeling Analysis Team*
\par
\end{document}